\section{Kết luận}
\label{sec:ketluan}

\subsection{Tổng kết kết quả thực hiện}
\label{sec:summary}

Đồ án đã xây dựng và vận hành hoàn chỉnh luồng xử lý dữ liệu GA4 từ khâu thu thập đến trực quan hóa: publisher phát lại các sự kiện theo đúng trình tự thời gian, Pub/Sub Emulator đóng vai trò hàng đợi trung gian cục bộ và Apache Beam xử lý dữ liệu thông qua DirectRunner. Sau khi chuẩn hóa, các tệp JSONL được nạp lên BigQuery Sandbox để dbt triển khai các tầng staging, mô hình sao (core star schema), data marts và các khung nhìn chỉ số (KPI views). Dashboard Streamlit khai thác trực tiếp các bảng tổng hợp này nhằm bảo đảm tính nhất quán thay vì tự tính toán một bộ chỉ số riêng lẻ.\medskip

Về mặt tổ chức dữ liệu, mô hình kho gồm 4 bảng fact dùng chung 8 bảng dimension với mức độ chi tiết (grain) được định nghĩa rõ ràng. Lớp tổng hợp chỉ lưu trữ các số đo có tính cộng được để tránh sai số khi tính toán lại các tỷ lệ ở nhiều chiều phân tích. Về mặt nghiệp vụ, mỗi phiên được xác định dựa trên cặp định danh người dùng và mã phiên, đồng thời gắn ngày theo sự kiện phát sinh đầu tiên. Giao dịch hợp lệ phải là sự kiện \texttt{purchase} có mã giao dịch thực, doanh thu dương và được khử trùng bằng cách giữ lại bản ghi sớm nhất. Khi phân bổ nguồn truy cập, hệ thống ưu tiên lấy nguồn cấp phiên và chỉ dùng nguồn cấp người dùng để thay thế nếu phiên không ghi nhận thông tin nguồn.\medskip

Kết quả kiểm thử thực tế cho thấy các thành phần trong kiến trúc đã phối hợp trơn tru: 100 sự kiện từ nguồn dữ liệu công khai được truyền qua emulator và DirectRunner, nạp đầy đủ vào tầng dữ liệu thô, dựng thành công 19 mô hình dbt, vượt qua toàn bộ 5 bài kiểm thử dữ liệu và hiển thị chính xác lên dashboard Streamlit. Toàn bộ luồng xử lý cũng đã được đóng gói và tái lập thành công trong môi trường Docker trên WSL, đi kèm cơ chế giám sát giới hạn tài nguyên và dọn dẹp tiến trình tự động. Quá trình vận hành trên BigQuery Sandbox luôn nằm trong hạn mức miễn phí nhờ áp dụng kiểm tra trước (dry-run), giới hạn dung lượng quét cho từng loại truy vấn và theo dõi mức sử dụng trên toàn dự án.\medskip

Tuy nhiên, các kết quả kiểm chứng hiện tại mới chỉ dừng ở mức nguyên mẫu (prototype) và chưa phản ánh đầy đủ hiệu năng trên toàn bộ tập dữ liệu lớn. Do lô dữ liệu thử nghiệm (mini-batch) không phát sinh giao dịch mua hàng hay thông tin mặt hàng, luồng xử lý doanh thu vẫn cần được kiểm thử thêm trên một lô dữ liệu riêng có kiểm soát. Bên cạnh đó, DirectRunner và Pub/Sub Emulator chỉ phục vụ mục đích kiểm thử cục bộ chứ chưa thay thế được hạ tầng lưu trữ và xử lý phân tán bền vững, còn tệp xuất CSV cho Looker Studio mới chỉ đáp ứng nhu cầu báo cáo tĩnh (snapshot) thay vì kết nối dữ liệu trực tiếp.\medskip

\subsection{Hướng mở rộng hệ thống}
\label{sec:extensions}

Các dịch vụ và giải pháp có liên quan đến tài khoản thanh toán (Billing Project) được tách riêng trong mục này do nằm ngoài phạm vi vận hành miễn phí hiện tại.\medskip

Khi cần mở rộng quy mô xử lý dữ liệu hoặc chuyển sang vận hành trên hạ tầng đám mây được quản lý toàn phần, hệ thống có thể nâng cấp tầng thu thập lên Cloud Pub/Sub, tầng xử lý lên Dataflow và sử dụng Cloud Storage làm vùng lưu trữ trung gian. Quá trình chuyển đổi này đòi hỏi chuẩn bị đầy đủ về hạn mức tài nguyên (quota), phân quyền IAM, cấu hình bucket, vòng đời tác vụ (job lifecycle) và quản lý chi phí, thay vì xem như các thành phần đã sẵn có trong bản nguyên mẫu. Song song đó, cần bổ sung mô hình dữ liệu cho chi phí quảng cáo và chi phí vận hành thực tế để có cơ sở tính toán chính xác các chỉ số như ROAS hay CPA.\medskip

Đối với nhu cầu báo cáo nâng cao, tài liệu kỹ thuật của Looker Studio chỉ ra rằng trình kết nối BigQuery trực tiếp yêu cầu một dự án có gắn tài khoản thanh toán cùng quyền đọc dữ liệu và thực thi truy vấn tương ứng \parencite{looker-bigquery}. Vì vậy, nếu cần cập nhật báo cáo theo thời gian thực, doanh nghiệp cần thiết lập dự án hợp lệ và kiểm soát chặt chẽ dung lượng quét của dashboard; phương án xuất tệp CSV từ Sandbox hiện tại chỉ là giải pháp tạm thời phù hợp với phạm vi miễn phí chứ chưa thể thay thế hoàn toàn kết nối tự động liên tục.\medskip

Cuối cùng, khi khối lượng dữ liệu và yêu cầu nghiệp vụ gia tăng, kiến trúc hệ thống cần tiếp tục hoàn thiện theo hướng nạp dữ liệu tăng dần (incremental load) thay cho nạp lại toàn bộ (full-refresh), thiết lập chính sách lưu trữ dài hạn, xây dựng cơ chế khôi phục khi bảng trên Sandbox hết hạn, cũng như kiểm chứng lại toàn bộ luồng doanh thu trên các lô dữ liệu có phát sinh giao dịch thực tế. Mọi bước mở rộng về hạ tầng đều cần đi kèm quy trình kiểm thử tương đương trên dữ liệu mới thay vì chỉ dựa vào kết quả của lô thử nghiệm nhỏ hiện tại.\medskip